% Compile with XeLaTeX (twice for cross-references), or Tectonic.
% Standalone CVPR-style course proposal: no cvpr.sty or .bib file required.
% This approximates CVPR typography; it is not an official submission template.
% Chinese names are retained exactly as supplied; all proposal prose is English.
\documentclass[10pt,twocolumn,letterpaper]{article}
\usepackage[letterpaper,textwidth=6.875in,textheight=8.875in,top=1in,columnsep=0.3125in]{geometry}
\usepackage{fontspec}
\setmainfont{texgyretermes}[
  Extension=.otf,
  UprightFont=*-regular,
  BoldFont=*-bold,
  ItalicFont=*-italic,
  BoldItalicFont=*-bolditalic]
\usepackage{xeCJK}
% TeX Live's AR PL Mingti2L Big5 font supports Traditional Chinese names.
\setCJKmainfont{bsmi00lp.ttf}
% Report missing glyphs as errors rather than silently omitting names.
\tracinglostchars=3
\usepackage{amsmath,amssymb,booktabs}
\usepackage{titlesec}
\usepackage[hidelinks]{hyperref}
\hypersetup{pdftitle={Shaolin Duck Soccer: Autonomous Microduck Shooting with Hierarchical Reinforcement Learning}}
\titleformat{\section}{\large\bfseries}{\thesection.}{0.5em}{}
\titleformat{\subsection}{\normalsize\bfseries}{\thesubsection.}{0.5em}{}
\titlespacing*{\section}{0pt}{9pt plus 2pt minus 1pt}{4pt}
\titlespacing*{\subsection}{0pt}{6pt plus 1pt minus 1pt}{3pt}
\setlength{\parindent}{1em}
\setlength{\parskip}{0pt}
\setlength{\emergencystretch}{1em}
\setlength{\abovedisplayskip}{5pt}
\setlength{\belowdisplayskip}{5pt}
\pagestyle{plain}
% Guard the assignment's two-page limit, including references.
\AtEndDocument{\ifnum\value{page}>2
  \PackageError{microduck-proposal}{The proposal exceeds two pages}
    {Shorten the text to meet the assignment's page limit.}
\fi}
\begin{document}
\twocolumn[{
\centering
{\fontsize{15}{17}\selectfont\bfseries Shaolin Duck Soccer: Autonomous Microduck Shooting\\
with Hierarchical Reinforcement Learning\par}
\vspace{9pt}
{\normalsize\bfseries Team Members\par}
\vspace{3pt}
{\normalsize
曾令城 (F74122030)\quad 黃威誠 (F74125313)\quad
葉沁蓉 (F74126238)\quad 李和穆 (E64112287)\par}
\vspace{3pt}
{\normalsize National Cheng Kung University\par}
\vspace{13pt}
}]

\begin{center}\bfseries Abstract\end{center}
\noindent\textit{We propose Shaolin Duck Soccer, a simulated bipedal robot that autonomously approaches a ball, aligns with a goal, and initiates a shot. A hierarchical controller will reuse fixed walking and kicking policies while a high-level policy, trained with Proximal Policy Optimization (PPO), selects movement commands and kick timing from simulator state. A curriculum will gradually expand initial distances, heading variations, and goal configurations. We will compare a geometric rule-based controller, high-level PPO without a curriculum, and PPO with a curriculum using identical low-level skills. Evaluation will measure valid goal rate, completion time, falls, and generalization. Deliverables include a reproducible simulation, autonomous shooting demonstrations, and a controlled analysis of skill composition and training schedules. Vision and hardware deployment are outside the core scope; two-robot passing is an optional extension.}

\section{Introduction}
Walking and kicking skills alone do not solve autonomous shooting. A robot must approach the ball from a suitable direction, establish a foot-specific stance, and switch skills without losing balance. Simply walking toward the ball can produce an inaccurate shot. Our central question is whether a learned high-level controller can coordinate existing skills more effectively across changing ball--goal configurations, and whether curriculum learning improves training efficiency.

Microduck provides an open platform for reusing locomotion and kicking policies~\cite{microduck}. Prior work studies transitions between primitive skills~\cite{lee}, agile bipedal soccer~\cite{haarnoja}, and directional striking under noisy sensing~\cite{xu}. SkillX investigates unified soccer skills~\cite{skillx}, while MARLadona addresses cooperative team play~\cite{marladona}. We target a smaller, reproducible task on Microduck rather than claim a new soccer capability: a single robot, an initially stationary ball, and an undefended goal on flat ground. Observations come from the simulator; camera perception, opponents, and physical deployment are excluded. The contribution will be the task implementation and a controlled comparison of decision strategies and training curricula.

\section{Method}
\subsection{Environment and Skill Reuse}
We will extend the official Microduck simulation with a goal, randomized episode initialization, termination conditions, and evaluation logging. Public ONNX policies will support initial skill validation. If parallel training requires compatible PyTorch checkpoints, we will obtain them or reproduce the required skills using the official training setup. All methods will share the same frozen low-level policies.

A valid goal requires the ball to cross the goal plane from the field into the opening between the posts. To distinguish shooting from continuous pushing, we will also record kick activation, kicking-foot contact, and subsequent ball separation. Falls, timeouts, and out-of-bounds events will terminate unsuccessful episodes.

\subsection{Hierarchical Controller}
The high-level policy receives robot proprioception, robot-relative ball position and velocity, relative goal center and orientation, and the current skill and its progress. Its initial action interface is
\begin{equation}
 a_t^H=(v_x^{\mathrm{cmd}},v_y^{\mathrm{cmd}},\omega_z^{\mathrm{cmd}},z_{\mathrm{kick}}).
\end{equation}
The first three components command walking; a fixed threshold on the fourth triggers a kick, rather than setting kick strength. Shot direction is primarily determined by stance and heading. We will initially update the high-level policy at 10\,Hz and the low-level controller at 50\,Hz. A triggered kick will remain active for a validated duration to prevent rapid skill switching.

The geometric baseline will approach a foot-specific preparation pose, align with the goal, slow down, and trigger the same kick skill. For ball position $p_b$ and goal center $p_g$, define
\begin{equation}
 u_g=\frac{p_g-p_b}{\|p_g-p_b\|},\qquad
 p_r^*=p_b-R(\theta_g)r_{\mathrm{kick}},
\end{equation}
where $\theta_g$ is the heading of $u_g$, and $r_{\mathrm{kick}}$ is the calibrated ball offset for the selected kicking foot. This provides a meaningful baseline that accounts for shooting geometry.

\subsection{PPO and Curriculum}
We will train the high-level actor and critic with PPO~\cite{ppo}, initially using separate two-layer MLPs with 128 units per layer and normalized observations. Rewards will combine a once-per-episode valid-goal bonus, progress toward the preparation pose, improved alignment, and goal-directed ball motion after a valid kick, with penalties for falls, leaving the field, and elapsed time. Progress-based or bounded shaping will limit repeated reward collection without shooting. Both PPO variants will use identical rewards.

The curriculum will progress from near-ready shooting poses, to nearby starts with heading offsets, and finally to wider starting distances and varied goal configurations. Tentative distance ranges are 0.2--0.5\,m and 0.5--1.0\,m for the latter stages, subject to feasibility checks. Difficulty may increase when recent training success reaches 70\%, while retaining some easier episodes. These are initial design choices, not established optimal settings.

\subsection{Experimental Protocol}
We will compare the following methods:
\begin{center}
\small
\begin{tabular}{@{}cll@{}}
\toprule
 & High-level controller & Curriculum \\
\midrule
A & Geometry + state machine & -- \\
B & PPO & No \\
C & PPO & Yes \\
\bottomrule
\end{tabular}
\end{center}
Methods B and C will share network size, rewards, and total environment interaction budgets, including all curriculum stages. Method B will sample the full training distribution from the outset. All three methods will use identical skills and held-out test configurations. We target at least three training seeds and 100 evaluation episodes per final policy per condition. Test configurations will be fixed before tuning. We plan to use two borrowed RTX 4090 GPUs for independent runs after single-GPU validation.

\section{Expected Results}
We expect to deliver a system demonstrating approach, alignment, and shooting from frontal, lateral, and reversed-heading starts, including changed goal locations. Achieving at least 70\% valid goals over 100 standard test episodes is an engineering acceptance target, not a measured result or guaranteed prediction.

The primary metric will be valid goal rate over \emph{all} episodes, including falls, timeouts, and missed contacts. We will report successful-episode completion time separately from timeout frequency, along with fall rate, post-contact shot-direction error, and goal rate under held-out poses and mild physical variations. Episodes without a valid kick will remain failures. Learning efficiency will be measured in environment interactions needed to reach a specified success rate.

Results will include goal-rate learning curves with variation across seeds, spatial success maps, and matched trajectory and failure-case comparisons. We will test whether curricula improve convergence speed or final performance without assuming PPO must outperform geometry. Failures caused by low-level skills will be distinguished from high-level decision errors. Deliverables will include code, models, configurations, evaluation procedures, and success/failure videos. Fixed-target passing and, subsequently, two robots sharing one physical simulation are stretch goals informed by cooperative soccer research~\cite{marladona}.

\begin{thebibliography}{7}
\small
\setlength{\itemsep}{1pt}
\setlength{\parskip}{0pt}
\bibitem{microduck}
Pollen Robotics. \href{https://github.com/pollen-robotics/microduck_rl}{Microduck RL: RL training environments for Microduck}. Open-source software, accessed Sept.\ 29, 2026.
\bibitem{ppo}
J.~Schulman, F.~Wolski, P.~Dhariwal, A.~Radford, and O.~Klimov. \href{https://arxiv.org/abs/1707.06347}{Proximal policy optimization algorithms}. arXiv:1707.06347, 2017.
\bibitem{lee}
Y.~Lee, S.-H.~Sun, S.~Somasundaram, E.~S.~Hu, and J.~J.~Lim. \href{https://youngwoon.github.io/transition/}{Composing complex skills by learning transition policies}. In \emph{ICLR}, 2019.
\bibitem{haarnoja}
T.~Haarnoja et~al. \href{https://doi.org/10.1126/scirobotics.adi8022}{Learning agile soccer skills for a bipedal robot with deep reinforcement learning}. \emph{Science Robotics}, 2024.
\bibitem{xu}
Z.~Xu et~al. \href{https://arxiv.org/abs/2512.06571}{Learning agile striker skills for humanoid soccer robots from noisy sensory input}. In \emph{ICRA}, 2026. arXiv:2512.06571.
\bibitem{skillx}
Z.~Ye et~al. \href{https://arxiv.org/abs/2609.06718}{SkillX: Unified multi-skill policy learning for humanoid soccer}. arXiv:2609.06718, 2026.
\bibitem{marladona}
Z.~Li, F.~Bjelonic, V.~Klemm, and M.~Hutter. \href{https://arxiv.org/abs/2409.20326}{MARLadona---Towards cooperative team play using multi-agent reinforcement learning}. arXiv:2409.20326, 2024; revised 2025.
\end{thebibliography}
\end{document}
